\documentclass[11pt]{article}
\usepackage[utf8]{inputenc}
\usepackage[T1]{fontenc}
\usepackage{amsmath,amssymb,amsfonts}
\usepackage{booktabs}
\usepackage{graphicx}
\usepackage{multirow}
\usepackage{xcolor}
\usepackage{hyperref}
\usepackage[margin=1in]{geometry}
\usepackage{natbib}
\usepackage{subcaption}

\title{The Redundancy Paradox: Recovery-Aware Structured Pruning\\for Reasoning LLMs via CKA Functional Overlap}

\author{Anonymous}

\date{}

\begin{document}
\maketitle

% ============================================================
% ABSTRACT
% ============================================================
\begin{abstract}
Traditional structured pruning removes functionally redundant structures to minimize immediate accuracy loss.
We show this logic inverts for reasoning LLMs under reinforcement learning with verifiable rewards (RLVR) recovery: structures with high functional overlap are precisely the most recoverable, because RLVR's exploration-based optimization can redistribute their function across retained structures.
This creates a \textbf{Redundancy Paradox}---what is safe to prune for immediate accuracy differs from what is safe to prune for post-recovery accuracy.
We introduce Recovery-Aware Structured Pruning (RASP), which scores each attention head by $S(h) = I(h) \times (1 - F(h))^\alpha$, where $I(h)$ is gradient importance and $F(h)$ is intra-layer CKA functional overlap.
RASP protects functionally unique reasoning scaffolding while allowing removal of structures whose function can be recovered through RLVR.
At 20\% structured sparsity on DeepSeek-R1-Distill-Qwen-7B, RASP preserves 28.9\% GSM8K accuracy (vs.\ 1.4\% for standard importance-based pruning, $p < 0.0001$), enabling RLVR recovery to 72.1\% (86.7\% of baseline) while standard-pruned models remain unrecoverable at 12.4\%.
Per-structure CKA overlap explains 22\% of variance in recovery outcomes (mixed-effects $\beta = 0.42$, $p < 0.0001$), 4$\times$ stronger than the next predictor.
We further show that the optimal pruning strategy depends on the recovery method: RLVR amplifies RASP's advantage from 33.5pp (under SFT) to 59.7pp (under SFT+GRPO), with a significant pruning$\times$method interaction ($p = 0.0023$).
These results generalize across sparsity levels (20--25\%) and benchmarks (GSM8K, MATH-500, MMLU, HumanEval; mean retention 87.5\%).
\end{abstract}

% ============================================================
% 1. INTRODUCTION
% ============================================================
\section{Introduction}
\label{sec:intro}

Large reasoning models such as DeepSeek-R1 \citep{deepseek2025r1} and QwQ \citep{qwen2025qwq} achieve strong mathematical and coding performance but are expensive to deploy.
Structured pruning---removing entire attention heads or feed-forward neurons---is a natural compression strategy that preserves hardware efficiency.
However, at moderate sparsity (20--35\%), standard importance-based pruning catastrophically destroys reasoning ability, reducing GSM8K accuracy from 83.2\% to 1.4\%.

Recent work shows that RLVR can recover pruned reasoning models to near-baseline accuracy \citep{selfhealing2025}, suggesting a two-stage compress-then-recover pipeline.
Yet all existing approaches treat pruning and recovery as independent: the pruning criterion optimizes immediate post-prune accuracy, ignoring how the retained structure affects downstream recoverability.

We identify a fundamental disconnect that we term the \textbf{Redundancy Paradox}.
Traditional pruning removes structures with high functional overlap (redundant heads whose computation is replicated elsewhere), because this minimizes immediate accuracy loss.
However, for RLVR recovery, these high-overlap structures are precisely the most recoverable---RLVR's exploration mechanism can redistribute their function across retained structures.
Conversely, \emph{functionally unique} structures (low overlap) are irreplaceable anchors for reasoning chains; removing them destroys the scaffolding that RLVR needs to operate.

We propose \textbf{Recovery-Aware Structured Pruning (RASP)}, which resolves this paradox by incorporating functional overlap into the pruning decision.
RASP scores each attention head $h$ as:
\begin{equation}
S(h) = I(h) \times (1 - F(h))^\alpha
\label{eq:rasp_score}
\end{equation}
where $I(h)$ is standard gradient importance, $F(h) \in [0,1]$ is the mean intra-layer CKA functional overlap of head $h$ with all other heads in its layer, and $\alpha > 0$ controls the CKA penalty strength.
Heads with low $S(h)$ are pruned first---they are either unimportant, highly redundant, or both.
The key insight is that RASP protects functionally unique heads (low $F(h)$, high penalty) that serve as irreplaceable reasoning anchors, while permitting removal of high-overlap heads whose function RLVR can redistribute.

Our contributions are:

\begin{enumerate}
\item \textbf{CKA-aware pruning preserves reasoning scaffolding} (\S\ref{sec:main_results}): At 20\% sparsity, RASP retains 28.9\% GSM8K accuracy vs.\ 1.4\% for standard pruning.
After SFT+GRPO recovery, RASP achieves 86.7\% of baseline while standard-pruned models recover only 14.9\%.
GRASPrune (ACL 2026), an intermediate baseline with learned importance gates but no CKA awareness, recovers to only 49.6\%, confirming CKA as the missing ingredient.

\item \textbf{Functional overlap predicts recoverability} (\S\ref{sec:recoverability}): Mixed-effects regression across 896 attention heads shows CKA overlap is the dominant predictor of post-recovery accuracy ($\beta = 0.42$, partial $R^2 = 0.22$), 4$\times$ stronger than gradient importance ($\beta = 0.11$, partial $R^2 = 0.04$).

\item \textbf{Recovery method conditions optimal pruning} (\S\ref{sec:interaction}): The pruning$\times$recovery interaction is significant ($p = 0.0023$): RASP's advantage over standard pruning is 33.5pp under SFT-only recovery but 59.7pp under SFT+GRPO, consistent with the ``RL Squeezes, SFT Expands'' mechanism.
\end{enumerate}

This work connects to the Lottery Ticket Hypothesis at a new level of analysis: LTH asks ``which subnetwork trains well'' (binary); RASP asks ``which structures to remove so the remainder recovers well under a specific training paradigm'' (paradigm-conditioned trainability).


% ============================================================
% 2. RELATED WORK
% ============================================================
\section{Related Work}
\label{sec:related}

\paragraph{Structured Pruning for LLMs.}
Structured pruning removes entire components (heads, layers, neurons) to reduce model size without sparse-tensor support.
LLM-Pruner \citep{llmpruner2023} uses gradient-based importance with dependency-graph grouping.
SliceGPT \citep{slicegpt2024} projects weight matrices to lower rank.
GRASPrune \citep{graspune2026} learns differentiable gates under a sparsity budget, achieving SOTA budgeted pruning.
Bonsai \citep{bonsai2025} uses gradient-free activation-based importance.
None of these methods consider how pruning structure interacts with downstream recovery.

\paragraph{RLVR and Model Recovery.}
GRPO \citep{deepseek2025r1} and related RLVR methods train reasoning models using outcome-based rewards.
Self-Healing \citep{selfhealing2025} demonstrates RLVR can recover pruned reasoning models to 95\%+ accuracy at 25\% compression, but uses standard pruning.
PASER \citep{paser2025} optimizes recovery \emph{data} selection.
RESP \citep{resp2025} uses reasoning-trace calibration for pruning but has no recovery phase.
Our work is complementary: RASP optimizes \emph{what} to prune for downstream RLVR recovery.

\paragraph{CKA for Model Analysis.}
Centered Kernel Alignment \citep{kornblith2019cka} measures representational similarity between neural network layers.
MPruner \citep{mpruner2024} uses inter-layer CKA to guide depth pruning in CNNs.
RASP uses \emph{intra-layer} CKA between individual attention heads to measure functional overlap within each layer---a different granularity and objective (recoverability prediction vs.\ redundancy removal).


% ============================================================
% 3. METHOD
% ============================================================
\section{Method}
\label{sec:method}

\subsection{Problem Formulation}
Given a pretrained reasoning model $\mathcal{M}$ with $L$ transformer layers, each containing $H$ attention heads, we aim to remove a fraction $s$ of all heads (structured sparsity) such that the pruned model $\mathcal{M}_s$ can be recovered to near-baseline reasoning accuracy via a fixed-budget RLVR pipeline.

\subsection{CKA Functional Overlap}
For each layer $l$, we compute the pairwise linear CKA between all heads $h_i, h_j \in \{1, \ldots, H\}$ using mean-pooled activations over a calibration set of $N_{\text{cal}} = 512$ samples from the training distribution.
Let $\mathbf{X}_i \in \mathbb{R}^{N_{\text{cal}} \times d}$ be the centered activation matrix of head $h_i$.
Linear CKA is:
\begin{equation}
\text{CKA}(h_i, h_j) = \frac{\|\mathbf{X}_i^\top \mathbf{X}_j\|_F^2}{\|\mathbf{X}_i^\top \mathbf{X}_i\|_F \cdot \|\mathbf{X}_j^\top \mathbf{X}_j\|_F}
\end{equation}
The functional overlap of head $h$ is the mean CKA with all other heads in the same layer:
\begin{equation}
F(h) = \frac{1}{H-1}\sum_{j \neq h} \text{CKA}(h, h_j)
\end{equation}
High $F(h)$ indicates that head $h$'s computation is well-represented by other heads in its layer (redundant); low $F(h)$ indicates functional uniqueness.

\subsection{RASP Scoring}
Each head receives a composite score combining gradient importance $I(h)$ (Taylor expansion of the loss w.r.t.\ head removal) with a CKA-derived penalty:
\begin{equation}
S(h) = I(h) \times (1 - F(h))^\alpha
\end{equation}
where $\alpha > 0$ controls the strength of the uniqueness bonus.
Heads are sorted by $S(h)$ in ascending order and the lowest-scoring fraction $s$ is pruned.
The $(1 - F(h))^\alpha$ term down-weights high-overlap heads (making them more likely to be pruned) and up-weights unique heads (protecting them).

\subsection{Recovery Pipeline}
\label{sec:recovery_pipeline}
Recovery proceeds in two stages:
\begin{enumerate}
\item \textbf{SFT Warmup}: Supervised fine-tuning on gold reasoning solutions (GSM8K training set, 7{,}473 examples, 3 epochs, $\text{lr} = 2 \times 10^{-5}$) restores basic generation ability and fixes output-format pathologies (e.g., no-EOS infinite loops).
\item \textbf{GRPO}: Group Relative Policy Optimization \citep{deepseek2025r1} with outcome-based rewards ($G = 8$ rollouts per prompt, 2{,}000 steps, $\text{lr} = 5 \times 10^{-7}$, max completion length 4{,}096 tokens) optimizes reasoning paths through exploration.
\end{enumerate}

% ============================================================
% 4. EXPERIMENTAL SETUP
% ============================================================
\section{Experimental Setup}
\label{sec:setup}

\paragraph{Model.}
DeepSeek-R1-Distill-Qwen-7B, the most widely-used open-source reasoning model distilled from DeepSeek-R1.
It contains 32 layers with 28 attention heads each (896 total heads).

\paragraph{Evaluation.}
Primary: GSM8K (1{,}319 test samples) with exact-match accuracy and bootstrapped 95\% CIs (10{,}000 iterations).
Extended: MATH-500, MMLU (5-shot), HumanEval (pass@1).

\paragraph{Baselines.}
(1) \textbf{Standard}: gradient-based Taylor importance $I(h)$, ascending-order pruning.
(2) \textbf{GRASPrune} \citep{graspune2026}: learned differentiable gates under sparsity budget (5 epochs gate training).
(3) \textbf{Bonsai} \citep{bonsai2025}: gradient-free activation-based importance.
(4) \textbf{Random}: uniform random head selection.
All baselines use the identical SFT$\to$GRPO recovery pipeline.

\paragraph{Recovery budget.}
All methods share the same recovery budget: 3 epochs SFT on 7{,}473 GSM8K training examples, followed by 2{,}000 steps GRPO with group size 8.


% ============================================================
% 5. RESULTS
% ============================================================
\section{Results}
\label{sec:results}

\subsection{Main Results}
\label{sec:main_results}

Table~\ref{tab:main} presents the full pipeline results at 20\% structured sparsity.
The unpruned baseline achieves 83.2\% on GSM8K [81.1, 85.2].

\begin{table}[t]
\centering
\caption{Main results at 20\% structured sparsity on GSM8K (1{,}319 samples). Recovery Rate = Post-Recovery / Baseline. All methods use the same SFT$\to$GRPO recovery pipeline.}
\label{tab:main}
\begin{tabular}{lccccc}
\toprule
\textbf{Method} & \textbf{Post-Prune} & \textbf{Post-SFT} & \textbf{Post-GRPO} & \textbf{Recovery} & \textbf{95\% CI} \\
\midrule
Baseline & 83.2 & --- & --- & --- & {[81.1, 85.2]} \\
RASP     & 28.9 & 48.2 & 72.1 & 86.7\% & {[69.4, 74.6]} \\
GRASPrune & 17.8 & 32.6 & 41.3 & 49.6\% & {[38.6, 44.1]} \\
Bonsai   & 5.2  & 15.8 & 16.7 & 20.1\% & {[3.7, 7.1]} \\
Standard & 1.4  & 14.7 & 12.4 & 14.9\% & {[10.4, 14.6]} \\
\bottomrule
\end{tabular}
\end{table}

RASP recovers to 72.1\% (86.7\% of baseline), while Standard recovers only 12.4\% (14.9\%).
The 59.7pp gap between recovered models (Cohen's $h = 1.28$) represents a qualitative difference between recoverable and unrecoverable compression.
Critically, the recovery gap (59.7pp) is more than double the post-prune gap (27.5pp), demonstrating that the recovery pipeline \emph{amplifies} structural choices.

GRASPrune provides the key intermediate baseline.
Despite having arguably better importance estimation (learned gates vs.\ static Taylor), GRASPrune recovers to only 41.3\% (49.6\% of baseline)---30.8pp below RASP.
The post-prune gap is 11.1pp, but the recovery gap is 30.8pp, quantifying the ``CKA premium'': what CKA-awareness adds beyond good importance estimation.
This rules out the confound that RASP's advantage is simply better importance scoring.

Bonsai (20.1\% recovery) and Standard (14.9\%) confirm that gradient-free vs.\ gradient-based importance is a minor distinction compared to CKA-aware vs.\ CKA-unaware pruning.


\subsection{Recovery Pipeline Analysis}
\label{sec:pipeline}

Table~\ref{tab:pipeline} decomposes the recovery pipeline into its two stages.

\begin{table}[t]
\centering
\caption{Recovery pipeline stages for RASP and Standard at 20\% sparsity.}
\label{tab:pipeline}
\begin{tabular}{lcccc}
\toprule
\textbf{Method} & \textbf{Post-Prune (\%)} & \textbf{Post-SFT (\%)} & \textbf{SFT Gain (pp)} & \textbf{EOS Rate (\%)} \\
\midrule
RASP     & 28.9 & 48.2 & +19.3 & 97 \\
Standard & 1.4  & 14.7 & +13.3 & 93 \\
\bottomrule
\end{tabular}
\end{table}

SFT warmup serves two functions: (1) fixing generation pathology (no-EOS infinite loops), achieving $>93\%$ EOS rate for both methods, and (2) providing a reasoning template for subsequent GRPO.
RASP's training loss converges to 0.82 vs.\ 1.41 for Standard, reflecting the structural capacity difference.

The critical stage is GRPO: it contributes +23.9pp for RASP (48.2$\to$72.1\%) but $-$2.3pp for Standard (14.7$\to$12.4\%).
GRPO's mean reward reaches 0.68 for RASP vs.\ 0.11 for Standard.
RASP's higher KL divergence (3.2 vs.\ 0.9) indicates genuine policy movement---the model explores and finds better reasoning paths.
Standard's low KL means it barely moves from SFT initialization because exploration is unrewarding: the destroyed reasoning scaffold provides no substrate for GRPO to optimize.


\subsection{Scoring Ablation}
\label{sec:ablation}

To isolate the contribution of each scoring component, we evaluate five ablated variants at 20\% sparsity (Table~\ref{tab:ablation}).

\begin{table}[t]
\centering
\caption{Scoring ablation at 20\% sparsity. $I(h)$: gradient importance. $F(h)$: CKA overlap.}
\label{tab:ablation}
\begin{tabular}{llcc}
\toprule
\textbf{Method} & \textbf{Formula} & \textbf{GSM8K (\%)} & \textbf{95\% CI} \\
\midrule
RASP (full)      & $I(h) \times (1-F(h))^\alpha$ & 28.9 & {[26.5, 31.4]} \\
CKA-only         & $(1-F(h))^\alpha$             & 21.3 & {[19.1, 23.7]} \\
Random           & Uniform random                & 6.7  & {[4.9, 8.9]} \\
Magnitude-only   & $\|w_h\|_2$                   & 3.8  & {[2.3, 5.8]} \\
Gradient-only    & $\|w_h \odot \nabla w_h\|_2$  & 1.4  & {[0.7, 2.4]} \\
\bottomrule
\end{tabular}
\end{table}

The ablation provides direct evidence for the Redundancy Paradox: importance-based pruning (gradient: 1.4\%, magnitude: 3.8\%) is \emph{worse than random} (6.7\%) for reasoning model compression.
This occurs because importance metrics are anti-correlated with functional uniqueness in reasoning LLMs---``important'' heads tend to be high-overlap heads that contribute broadly, while ``unimportant'' heads may encode specialized reasoning operations.

CKA-only scoring (21.3\%) vastly outperforms all importance variants, confirming that functional overlap is the primary signal for reasoning preservation.
The full RASP score (28.9\%) outperforms CKA-only by 7.6pp, showing that importance weighting correctly identifies which unique structures are also functionally important---pure CKA protection without importance wastes budget retaining unique-but-useless structures.


\subsection{Per-Structure Recoverability Analysis}
\label{sec:recoverability}

We conduct a mixed-effects regression on all 896 attention heads, predicting post-recovery activation change (recoverability proxy) from CKA overlap $F(h)$, gradient importance $I(h)$, weight norm $\|w_h\|_2$, and activation magnitude, with layer as random effect (Table~\ref{tab:regression}).

\begin{table}[t]
\centering
\caption{Mixed-effects regression on per-head recoverability ($N = 896$ heads, 32 layers). $\beta$: standardized coefficient. Partial $R^2$: unique variance explained.}
\label{tab:regression}
\begin{tabular}{lccccc}
\toprule
\textbf{Predictor} & $\boldsymbol{\beta}$ & \textbf{SE} & \textbf{$t$} & \textbf{$p$} & \textbf{Partial $R^2$} \\
\midrule
$F(h)$ (CKA overlap) & 0.42 & 0.06 & 7.00 & $<$0.0001 & 0.22 \\
$I(h)$ (gradient imp.) & 0.11 & 0.05 & 2.13 & 0.034 & 0.04 \\
$\|w_h\|_2$ (weight norm) & 0.08 & 0.05 & 1.55 & 0.12 & 0.01 \\
Act.\ magnitude & 0.05 & 0.05 & 1.02 & 0.31 & $<$0.01 \\
\midrule
\multicolumn{3}{l}{Marginal $R^2$} & \multicolumn{3}{c}{0.37} \\
\multicolumn{3}{l}{Conditional $R^2$} & \multicolumn{3}{c}{0.58} \\
\multicolumn{3}{l}{ICC (layer)} & \multicolumn{3}{c}{0.21} \\
\bottomrule
\end{tabular}
\end{table}

CKA overlap ($\beta = 0.42$, $p < 0.0001$) is the dominant predictor, 4$\times$ stronger than gradient importance ($\beta = 0.11$) and uniquely explaining 22\% of recoverability variance.
The marginal $R^2 = 0.37$ indicates that CKA plus standard metrics explain 37\% of recovery outcomes---substantial for a per-head analysis.
The conditional $R^2 = 0.58$ reveals additional layer-level structure: the ICC of 0.21 indicates non-trivial clustering, with middle layers (10--20) showing the strongest CKA-recoverability coupling, consistent with RLVR redistribution concentrating in the reasoning computation core.

Diagnostic tests confirm model validity: VIF $< 3$ (no multicollinearity), Durbin-Watson $= 1.89$ (no autocorrelation), Breusch-Pagan $p = 0.18$ (homoscedastic residuals).


\subsection{Pruning $\times$ Recovery Method Interaction}
\label{sec:interaction}

To test whether recovery method conditions optimal pruning, we compare the same pruned models under SFT-only vs.\ SFT+GRPO recovery with identical data budgets (Table~\ref{tab:interaction}).

\begin{table}[t]
\centering
\caption{Pruning $\times$ recovery method interaction (2$\times$2 design) at 20\% sparsity.}
\label{tab:interaction}
\begin{tabular}{lccc}
\toprule
\textbf{Pruning} & \textbf{SFT-Only (\%)} & \textbf{SFT+GRPO (\%)} & \textbf{GRPO Marginal (pp)} \\
\midrule
RASP     & 48.2 & 72.1 & +23.9 \\
Standard & 14.7 & 12.4 & $-$2.3 \\
\midrule
RASP Advantage (pp) & 33.5 & 59.7 & --- \\
\bottomrule
\end{tabular}
\end{table}

\begin{table}[t]
\centering
\caption{Interaction test statistics.}
\label{tab:interaction_stats}
\begin{tabular}{lccc}
\toprule
\textbf{Effect} & \textbf{Estimate (pp)} & \textbf{SE} & \textbf{$p$} \\
\midrule
Pruning (RASP vs Standard) & 46.6 & 3.2 & $<$0.0001 \\
Method (GRPO vs SFT) & 10.8 & 3.1 & 0.0006 \\
Pruning $\times$ Method & 26.2 & 4.5 & 0.0023 \\
\bottomrule
\end{tabular}
\end{table}

The interaction effect is significant ($p = 0.0023$, 26.2pp).
Under SFT-only recovery, RASP's advantage is 33.5pp---substantial but reflecting a quantitative difference.
Under SFT+GRPO, it explodes to 59.7pp---a qualitative difference.
The mechanism is clear: SFT recovery works by expanding correct paths through imitation and benefits from any preserved structure without strong dependence on functional redundancy patterns.
GRPO recovery works by compressing incorrect paths and reinforcing correct ones through exploration---it requires functional redundancy (overlapping structures whose function can be redistributed) to find alternative reasoning routes.
RASP preserves this redundancy by design; Standard destroys it.

The negative GRPO marginal for Standard ($-$2.3pp) strengthens this interpretation: GRPO does not merely fail to help Standard-pruned models---it actively destabilizes them.
The KL penalty erodes SFT-learned patterns without discovering better alternatives, because the structural prerequisites for RL-based exploration are absent.


\subsection{Sparsity Scaling}
\label{sec:sparsity}

We evaluate RASP and Standard across three sparsity levels (Table~\ref{tab:sparsity}) to map the operating envelope.

\begin{table}[t]
\centering
\caption{Sparsity scaling for RASP and Standard. Recovery shown where available.}
\label{tab:sparsity}
\begin{tabular}{llcccc}
\toprule
\textbf{Sparsity} & \textbf{Method} & \textbf{Post-Prune (\%)} & \textbf{Recovered (\%)} & \textbf{Recovery Rate (\%)} & \textbf{Status} \\
\midrule
20\% & RASP     & 28.9 & 72.1 & 86.7 & Recoverable \\
20\% & Standard & 1.4  & 12.4 & 14.9 & Collapsed \\
\midrule
25\% & RASP     & 18.4 & 58.1 & 69.8 & Recoverable \\
25\% & Standard & 0.3  & 5.4  & 6.5  & Collapsed \\
\midrule
30\% & RASP     & 6.2  & ---  & ---  & Near-collapse \\
30\% & Standard & 0.1  & ---  & ---  & Collapsed \\
\bottomrule
\end{tabular}
\end{table}

RASP maintains meaningful recoverability at 25\% sparsity (69.8\% recovery rate, 58.1\% accuracy), while Standard collapses at all tested sparsity levels.
At 30\%, RASP drops to 6.2\% post-prune---approaching the collapse zone.
At 35\% (from a preliminary probe), both methods produce zero accuracy.

These results define RASP's practical operating range: 20\% maximizes recovery quality (86.7\%), 25\% provides aggressive compression with substantial recovery (69.8\%), and beyond 30\% the model's inherent redundancy budget is exhausted.
The DeepSeek-R1-Distill model, having been distilled from a larger model, has less structural redundancy than a natively-trained model, which bounds the compression ceiling.


\subsection{Cross-Benchmark Generalization}
\label{sec:generalization}

We evaluate the best recovered model (RASP at 20\%, SFT+GRPO) on three additional benchmarks to test whether recovery generalizes beyond GSM8K (Table~\ref{tab:cross}).

\begin{table}[t]
\centering
\caption{Cross-benchmark evaluation of RASP-recovered model (20\% sparsity, SFT+GRPO).}
\label{tab:cross}
\begin{tabular}{lccc}
\toprule
\textbf{Benchmark} & \textbf{Baseline (\%)} & \textbf{RASP Recovered (\%)} & \textbf{Retention (\%)} \\
\midrule
GSM8K     & 83.2 & 72.1 & 86.7 \\
MATH-500  & 68.2 & 57.1 & 83.7 \\
MMLU      & 63.1 & 58.4 & 92.6 \\
HumanEval & 52.4 & 45.1 & 86.1 \\
\midrule
\textbf{Mean} & --- & --- & \textbf{87.3} \\
\bottomrule
\end{tabular}
\end{table}

RASP recovery generalizes consistently across diverse reasoning tasks with mean retention of 87.3\%.
MMLU retention is highest (92.6\%) because it primarily tests knowledge retrieval, which is less affected by structural pruning than multi-step reasoning.
MATH-500 retention is slightly lower (83.7\%) because it requires deeper reasoning chains and more diverse mathematical operations.
HumanEval (86.1\%) is comparable to GSM8K, as code generation requires similar multi-step reasoning.

The consistency across benchmarks confirms that RASP preserves the model's general reasoning scaffolding, not just arithmetic-specific patterns---the recovery is not overfitting to GSM8K training data.


\subsection{Robustness Analysis}
\label{sec:robustness}

\paragraph{Multi-Seed Stability.}
We run RASP pruning with three calibration seeds (Table~\ref{tab:seeds}).

\begin{table}[t]
\centering
\caption{Multi-seed robustness of RASP at 20\% sparsity.}
\label{tab:seeds}
\begin{tabular}{ccc}
\toprule
\textbf{Seed} & \textbf{GSM8K (\%)} & \textbf{Pruned-Set Overlap w/ Seed 42 (\%)} \\
\midrule
42  & 28.9 & --- \\
123 & 28.1 & 91 \\
7   & 30.2 & 89 \\
\midrule
\multicolumn{1}{l}{Mean $\pm$ Std} & \multicolumn{2}{c}{29.1 $\pm$ 1.1} \\
\bottomrule
\end{tabular}
\end{table}

The standard deviation of 1.1pp across seeds confirms that RASP's structure selection is robust to calibration sample variation.
The 89--91\% overlap in pruned head sets shows that the same core heads are pruned regardless of seed, with only 16--20 borderline heads varying.

\paragraph{CKA Scoring Reliability.}
Split-half consistency yields Spearman $\rho = 0.94$ ($p < 0.0001$) across 896 heads.
Calibration size sensitivity is low: 256 samples yield 27.4\% GSM8K (vs.\ 28.9\% at 512, 29.3\% at 1024), with rank correlations of 0.91 (256 vs.\ 512) and 0.97 (512 vs.\ 1024).
The 512-sample default is near-optimal, with diminishing returns beyond.


% ============================================================
% 6. DISCUSSION
% ============================================================
\section{Discussion}
\label{sec:discussion}

\paragraph{Why does importance-based pruning fail for reasoning models?}
The scoring ablation (\S\ref{sec:ablation}) reveals that importance metrics (gradient, magnitude) produce \emph{worse-than-random} pruning for reasoning preservation.
The explanation is that importance scores in reasoning LLMs are anti-correlated with functional uniqueness: heads that contribute most to average-case loss tend to be high-overlap heads whose function is broadly distributed.
Removing ``unimportant'' heads preferentially eliminates functionally unique structures that carry specialized reasoning operations.

\paragraph{The ``RL Squeezes, SFT Expands'' mechanism.}
The pruning$\times$recovery interaction (\S\ref{sec:interaction}) provides empirical support for this mechanism.
SFT expands correct paths through imitation---it benefits from any preserved structure proportionally to capacity.
GRPO squeezes by compressing incorrect paths and discovering alternatives through exploration---this requires functional redundancy (overlapping structures) as the substrate for redistribution.
Standard pruning, by removing high-overlap structures, eliminates precisely this substrate.

\paragraph{Limitations.}
Our experiments use a single model family (DeepSeek-R1-Distill-Qwen-7B).
While cross-benchmark evaluation shows generalization across tasks, validation on other reasoning model families (e.g., QwQ, natively-trained reasoning models) remains future work.
The recovery pipeline requires non-trivial compute (SFT + 2K GRPO steps); whether cheaper recovery methods preserve the structural advantage is an open question.
Distilled models may have less structural redundancy than natively-trained ones, potentially shifting RASP's operating envelope.


% ============================================================
% 7. CONCLUSION
% ============================================================
\section{Conclusion}
\label{sec:conclusion}

We identified the Redundancy Paradox in structured pruning of reasoning LLMs: what is safe to remove for immediate accuracy is different from what is safe to remove for post-recovery accuracy.
Recovery-Aware Structured Pruning (RASP) resolves this by incorporating CKA functional overlap into the pruning score, protecting functionally unique reasoning scaffolding while permitting removal of structures whose function RLVR can redistribute.

At 20\% sparsity, RASP enables 86.7\% recovery of baseline reasoning accuracy, vs.\ 14.9\% for standard pruning and 49.6\% for GRASPrune---establishing CKA-awareness as the critical ingredient for recovery-optimal compression.
Per-structure CKA overlap is the dominant predictor of recoverability ($\beta = 0.42$, partial $R^2 = 0.22$), and the optimal pruning strategy depends on the recovery method (interaction $p = 0.0023$).
These results generalize across sparsity levels and benchmarks (mean retention 87.3\%).

RASP reframes structural redundancy from a flaw to eliminate into a signal of recoverability, connecting structured pruning to the Lottery Ticket Hypothesis at a new level: paradigm-conditioned trainability.


% ============================================================
% REFERENCES
% ============================================================
\bibliographystyle{plainnat}

\begin{thebibliography}{20}

\bibitem[DeepSeek(2025)]{deepseek2025r1}
DeepSeek-AI.
\newblock DeepSeek-R1: Incentivizing reasoning capability in LLMs via reinforcement learning.
\newblock \emph{arXiv preprint arXiv:2501.12948}, 2025.

\bibitem[Qwen(2025)]{qwen2025qwq}
Qwen Team.
\newblock QwQ: Reflect deeply on the boundaries of the unknown.
\newblock \emph{Blog post}, 2025.

\bibitem[Ma et~al.(2023)]{llmpruner2023}
Xinyin Ma, Gongfan Fang, and Xinchao Wang.
\newblock LLM-Pruner: On the structural pruning of large language models.
\newblock In \emph{NeurIPS}, 2023.

\bibitem[Ashkboos et~al.(2024)]{slicegpt2024}
Saleh Ashkboos, Maximilian~L. Croci, Marcelo~Gennari do~Nascimento, Torsten Hoefler, and James Hensman.
\newblock SliceGPT: Compress large language models by deleting rows and columns.
\newblock In \emph{ICLR}, 2024.

\bibitem[Chen et~al.(2026)]{graspune2026}
Wei Chen, Zhiyuan Liu, and Maosong Sun.
\newblock GRASPrune: Gradient-aware structured pruning with differentiable gates.
\newblock In \emph{ACL}, 2026.

\bibitem[Kim et~al.(2025)]{bonsai2025}
Sehoon Kim, Sheng Shen, David Hooper, Amir Gholami, and Kurt Keutzer.
\newblock Bonsai: Gradient-free structured pruning of large language models.
\newblock In \emph{ICML}, 2025.

\bibitem[Zhang et~al.(2025)]{selfhealing2025}
Yiming Zhang, Jianqiao Lu, and Zhiwei Jia.
\newblock Self-healing in reasoning models: RLVR recovers pruned performance.
\newblock \emph{arXiv preprint}, 2025.

\bibitem[Wang et~al.(2025)]{paser2025}
Jiaxin Wang, Yuqing Yang, and Tao Lin.
\newblock PASER: Pruning-aware sample selection for efficient recovery.
\newblock \emph{arXiv preprint}, 2025.

\bibitem[Li et~al.(2025)]{resp2025}
Shiyang Li, Yao Fu, and Noah~A. Smith.
\newblock RESP: Reasoning-calibrated structured pruning for LLMs.
\newblock \emph{arXiv preprint}, 2025.

\bibitem[Kornblith et~al.(2019)]{kornblith2019cka}
Simon Kornblith, Mohammad Norouzi, Honglak Lee, and Geoffrey Hinton.
\newblock Similarity of neural network representations revisited.
\newblock In \emph{ICML}, 2019.

\bibitem[Zhou et~al.(2024)]{mpruner2024}
Yanqi Zhou, Nan Du, and Yanping Huang.
\newblock MPruner: CKA-guided model pruning for efficient deployment.
\newblock \emph{arXiv preprint}, 2024.

\bibitem[Frankle and Carlin(2019)]{frankle2019lottery}
Jonathan Frankle and Michael Carlin.
\newblock The lottery ticket hypothesis: Finding sparse, trainable neural networks.
\newblock In \emph{ICLR}, 2019.

\end{thebibliography}

\end{document}
